\documentclass[../BTL.tex]{subfiles}
\begin{document}

\section{Tài liệu tham khảo}

Trong quá trình thực hiện đề tài, nhóm đã tham khảo và trích dẫn các nguồn tài liệu khoa học uy tín từ các hội nghị và tạp chí hàng đầu trong lĩnh vực Deep Learning và Xử lý Ngôn ngữ Tự nhiên. Dưới đây là danh sách các tài liệu tham khảo chính được phân loại theo chủ đề.

Các tài liệu về kiến trúc mạng nơ-ron hồi quy (RNN) và các biến thể của nó bao gồm: bài báo gốc của Hochreiter và Schmidhuber năm 1997 giới thiệu kiến trúc Long Short-Term Memory (LSTM) được công bố trên tạp chí Neural Computation~\cite{hochreiter1997lstm}; bài báo của Cho và cộng sự năm 2014 giới thiệu Gated Recurrent Unit (GRU) tại hội nghị EMNLP~\cite{cho2014gru}; nghiên cứu của Graves và cộng sự về ứng dụng RNN trong nhận dạng giọng nói~\cite{graves2012supervised}; và tổng quan của Schmidhuber về Deep Learning trong mạng nơ-ron~\cite{schmidhuber2015deep}.

Các tài liệu về cơ chế Attention bao gồm: bài báo nền tảng của Bahdanau và cộng sự năm 2014 về Neural Machine Translation sử dụng cơ chế Attention và Align~\cite{bahdanau2014attention}; nghiên cứu của Yang và cộng sự năm 2016 về Hierarchical Attention Networks cho phân loại văn bản~\cite{yang2016hierarchical}; và bài báo của Zhou và Mao năm 2015 về mô hình C-LSTM cho phân loại văn bản~\cite{zhou2015attentionrcnn}.

Kiến trúc Transformer và các mô hình Pretrained được đề cập trong: bài báo "Attention Is All You Need" của Vaswani và cộng sự năm 2017 tại NeurIPS~\cite{vaswani2017transformer}; BERT của Devlin và cộng sự năm 2019 tại NAACL-HLT~\cite{devlin2019bert}; DistilBERT của Sanh và cộng sự năm 2019~\cite{sanh2019distilbert}; và tài liệu thư viện Transformers của Wolf và cộng sự năm 2020 tại EMNLP~\cite{wolf2020transformers}.

Các tài liệu về tối ưu hóa và huấn luyện bao gồm: Adam optimizer của Kingma và Ba năm 2015~\cite{kingma2015adam}; AdamW của Loshchilov và Hutter năm 2017 về Decoupled Weight Decay~\cite{loshchilov2017adamw}; Dropout của Srivastava và cộng sự năm 2014~\cite{srivastava2014dropout}; Batch Normalization của Ioffe và Szegedy năm 2015~\cite{ioffe2015batch}; Layer Normalization của Ba và cộng sự năm 2016~\cite{ba2016layer}; Focal Loss của Lin và cộng sự năm 2017 cho đối tượng Detection~\cite{lin2017focal}; và các nghiên cứu về khởi tạo trọng số của Glorot và Bengio~\cite{glorot2010deep} và He và cộng sự~\cite{he2015delving}.

Các tài liệu về độ đo đánh giá bao gồm: giới thiệu về ROC Analysis của Fawcett năm 2006~\cite{fawcett2006introduction}; phân tích hệ thống về các độ đo cho bài toán phân loại của Sokolova và Lapalme năm 2009~\cite{sokolova2009systematic}; và hướng dẫn đánh giá và lựa chọn mô hình của Raschka năm 2018~\cite{raschka2018roc}.

Các mô hình nhúng từ (Word Embeddings) được tham khảo từ: GloVe của Pennington và cộng sự năm 2014~\cite{pennington2014glove}; Word2Vec của Mikolov và cộng sự năm 2013~\cite{mikolov2013word2vec}; và CNN cho phân loại câu của Kim năm 2014~\cite{kim2014convolutional}.

Bộ dữ liệu Jigsaw Toxic Comment Classification Challenge được tham khảo từ cuộc thi chính thức trên Kaggle~\cite{kaggleJigsaw}. Các nghiên cứu liên quan đến hate speech detection bao gồm: Djuric và cộng sự năm 2015 về phát hiện harassment sử dụng comment embeddings~\cite{djuric2015hate}; và Warner và cộng sự năm 2012 về phát hiện harassment trên Twitter~\cite{warner2012detecting}.

Các tài liệu về RCNN và các kiến trúc kết hợp bao gồm: Lai và cộng sự năm 2015 về Recurrent Convolutional Neural Networks cho phân loại văn bản tại AAAI~\cite{lai2015rcnn}; và Socher và cộng sự năm 2013 về Recursive Deep Models for Sentiment Treebank tại EMNLP~\cite{socher2013recursive}.

Cuối cùng, các tài liệu tổng quát về Deep Learning bao gồm: cuốn sách "Deep Learning" của Goodfellow và cộng sự năm 2016~\cite{goodfellow2016deep}; "Neural Networks: Tricks of the Trade" của LeCun và cộng sự năm 1998 về các thủ thuật huấn luyện~\cite{lecun1998gradient}; bài báo về "No Free Lunch Theorems" của Wolpert năm 1997~\cite{wolpert1997no}; và tài liệu PyTorch của Paszke và cộng sự năm 2019 tại NeurIPS~\cite{paszke2019pytorch}.

\end{document}
